%قبل از شروع مسئله ترجیحاً چیزی ننویسید. اما اگر خواستید بنویسید، در قالب لیستی (فقرات و شمارش) بنویسید.
%برای اضافه کردن نمره، بعدی از شروع‌مسئله، [x] را اضافه کنید و به جای x  نمره مورد نظر را قرار دهید. برای زیرمسئله‌ها هم به همین صورت است.
%\شروع‌مسئله[10]
\شروع‌مسئله
% صورت اصلی مسئله را در این‌جا بنویسید.
(۳۰ نمره) برای هر یک از مسائل زیر الگوریتمی غیر بازگشتی با مرتبه‌زمانی خطی ارائه دهید.





%اگر مسئله چند قسمت داشت به این صورت عمل کنید:
\شروع{زیرمسائل}
\زیرمسئله
(۱۵ نمره) فرض کنید تنها حق ذخیره کردن دو عدد اضافه را در حافظه دارید. در اینصورت با ورودی گرفتن یک درخت دودویی جست‌وجو پیمایش پس‌ترتیب آن‌را ارائه کنید.
\زیرمسئله
(۱۵ نمره)
پیمایش پیش‌ترتیب یک درخت دودویی جست‌وجو با $n$ عضو داده‌شده است. با استفاده از یک پشته و بدون ساختن خود درخت برگ‌های آن را پیدا کنید.
\پایان{زیرمسائل}

\solution{
	\begin{enumerate}
		\item
		کافی است محل قبلی اشاره‌گر را در یک خانه از حافظه نگه‌داریم. با توجه به محل فعلی و محل قبلی، می‌توان حرکت بعدی را به دست آورد. فرض کنید اکنون اشاره‌گر روی گره‌ی X باشد و مرحله‌ی قبل روی پدر X بوده باشد. این یعنی تازه وارد این زیردرخت شده و باید به زیر درخت چپ برود. فرض کنید در حالتی دیگر، در مرحله‌ی قبلی روی فرزند چپ X بوده باشد. پس اکنون باید به فرزند راست X برود و اگر در مرحله‌ی قبلی روی فرزند راست X بوده است، یعنی این پیمایش این زیردرخت  تمام شده و خود X باید چاپ شود.
		\item
	لیست پیش‌ترتیب درخت با دو اشاره‌گر i و j پیمایش می‌شود که در ابتدا به خانه‌ی اول و دوم لیست اشاره می‌کنند. مقادیر منتاظر با هر کدام را در لیست، 
$p_1$
 و 
$p_2$
 می‌نامیم. تا جایی که 
$p_1>p_2$
 باشد، به این معناست که در زیردرخت چپ هستیم و هنوز برگی ملاقات نشده است. 
$p_1$
 را در پشته push می‌کنیم و هر دو  اشاره‌گر را یکی به جلو می‌بریم.

اگر
$p_2>p_1$
 شد و
$p_2$
 از مقدار بالای پشته بزرگتر بود، نشانه‌ی آن است که 
$p_1$
 یک برگ درخت است. تا جایی که
$p_2$
 از مقدار بالای پشته بزرگتر است، از پشته pop می‌کنیم. این روند تکرار شده و لیست یک بار پیمایش می‌شود. برگ‌ها در زمان خطی مشخص شده‌اند.
	\end{enumerate}
}




%%%%%%%%%%%%%%%%%%%%%%%%%%%%%%%%%%%%%%%%%%%%%%%%%%%%%
% برخی دستورات مفید برای تمرین‌ها. %
%%%%%%%%%%%%%%%%%%%%%%%%%%%%%%%%%%%%%%%%%%%%%%%%%%%%%

% فرمول‌های نمونه:
%$f(x) = \Theta(x)$
%$T(x, y) \leq \log(x) + 2^y + \binom{2n}{n}$
% {\tt code\_function}


% لیست موارد بدون شماره‌گذاری:
%\شروع{فقرات}
%    \فقره مورد اول 
%        \شروع{فقرات}
%            \فقره مورد اول داخلی 
%                \شروع{فقرات}
%                    \فقره اولین مورد خیلی داخلی 
%                    \فقره دومین مورد خیلی داخلی 
%                \پایان{فقرات}
%            \فقره مورد دوم داخلی 
%        \پایان{فقرات}
%    \فقره مورد دوم 
%\پایان{فقرات}

% لیست موارد شماره‌گذاری شده:
%\شروع{شمارش}
%    \فقره یک مورد 
%    \فقره یک مورد دیگر 
%    \فقره باز هم یک مورد دیگر
%\پایان{شمارش}

% برای نوشتن یک فرمول چند خطی که یک شماره بخوره. برای حالت‌های دیگر به سایت زیر مراجعه کنید:
% https://www.overleaf.com/learn/latex/Aligning_equations_with_amsmath
%\begin{align} 
%    \begin{split}
%        (x+y)^3 &= (x+y)^2(x+y) \\
%                &= (x^2+2xy+y^2)(x+y) \\
%                &= (x^3+2x^2y+xy^2) + (x^2y+2xy^2+y^3) \\
%                &= x^3+3x^2y+3xy^2+y^3
%    \end{split}                                 
%\end{align}

% نوشتن ماتریس و آرایه دوبعدی
%\begin{align}
%    A = 
%    \begin{bmatrix}
%        A_{11} & A_{21} \\
%        A_{21} & A_{22}
%    \end{bmatrix}
%\end{align}

% اضافه کردن تصویر با شرح. برای اضافه کردن تصویر بدون شرح، فقط دستور درج‌تصویر را استفاده کنید.
%\شروع{شکل}[h!]
%\تنظیم‌ازوسط
%  \درج‌تصویر[پهنا=0.5\پهنای‌متن]{img.jpg}
%\شرح{شرح تصویر}
%\برچسب{تصویر:۱}
%\پایان{شکل}
%\پایان{زیرمسائل}

%کد زیر را ببینید.
%%\lstinputlisting{source_filename.py}
%\begin{latin}
%\begin{lstlisting}
%# Import the modules
%import sys
%import random
%
%ans = True
%
%while ans:
%question = raw_input("Ask the magic 8 ball a question: (press enter to quit) ")
%
%answers = random.randint(1,8)
%
%if question == "":
%sys.exit()
%
%elif answers == 1:
%print "It is certain"
%
%elif answers == 2:
%print "Outlook good"
%
%elif answers == 3:
%print "You may rely on it"
%
%elif answers == 4:
%print "Ask again later"
%
%elif answers == 5:
%print "Concentrate and ask again"
%
%elif answers == 6:
%print "Reply hazy, try again"
%
%elif answers == 7:
%print "My reply is no"
%
%elif answers == 8:
%print "My sources say no"
%\end{lstlisting}
%\end{latin}
